% Page 037 — page imprimée 33
% Contenu seul : le préambule et la pagination sont gérés par meyrueis.tex



La caverne avait l’aspect d’une longue salle rectangulaire d’environ 20m sur 5, recouverte par
un rocher sans fissure. Nous marchions sur une couche élastique de très anciennes crottes de
mouton, parfaitement sèches et sans odeur. Au fond s’élevait une masse calcaire, un peu
comme un autel dans le chœur d’une église.
On pouvait en faire le tour par une galerie relativement étroite dans laquelle s’ouvraient
quelques trous à différentes hauteurs. L’un d’eux, au ras du sol, était plus profond que les
autres. C’était un véritable boyau dans lequel nous nous \textit{engageâmes d’abord en courbant le dos,
puis à genoux, enfin en rampant à mesure que la voûte s’abaissait au dessus de nous.} Le sol
était meuble, friable : un limon desséché depuis des siècles !

Frank était devant moi avec une lampe électrique. Après quelques minutes il s’arrêta un
moment, puis commença à reculer, lorsque nous nous relevâmes dans la galerie il était excité :
\textit{« j’ai peut-être trouvé ce que nous cherchions, me dit-il : Au fond du tunnel il y a un trou.
Pas bien grand mais je pouvais y passer le bras. Et de l’autre côté ma main ne rencontrait
plus de paroi, plus de sol, seulement un courant d’air. C’est sûrement une salle, peut-être
très grande, peut-être le réservoir de la Douce. Nous reviendrons demain. »}

Frank avait déjà une certaine expérience en spéléologie. Il connaissait la grotte de Dargilan, la
grande attraction touristique de la région. L’Aven Armand n’était pas encore exploité. Dans
cet immense labyrinthe il connaissait non seulement l’itinéraire ouvert au public (1h1/2 de
marche) mais encore salles et galeries restées secrètes. A l’entendre les plus belles,
naturellement ! Il savait donc ce qu’il nous fallait et lorsque nous entrâmes dans la Bergerie le
lendemain matin, nous avions un équipement rudimentaire mais suffisant : la forte lampe à
acétylène d’un vélo, une lampe électrique, 3 bougies, 2 boites d’allumettes et 3 pelotons de fil
rouge « empruntés » à la filature du Moulin d’Ayres, pour nous servir de fil d’Ariane… et une
solide pelle à charbon.

Le déblaiement du tunnel que nous pûmes approfondir d’environ 30 cm nous prit une bonne
heure, mais après ce travail nous pouvions nous y glisser à l’aise et surtout, la paroi terminale
n’était plus un obstacle : le trou était devenu une large ouverture. La lampe à acétylène nous
révéla que nous nous trouvions en haut d’une vaste salle très profonde, au sommet d’un
éboulis formé par des débris de stalagmites et stalactites éboulés sur le sol et sur lesquels, par
endroits, de nouvelles et importantes stalagmites s’étaient formées. Le même phénomène a été
observé dans l’Aven Armand et les géologues ont estimé que ces éboulis étaient dus à un
tremblement de terre ayant eu lieu il y a environ 15 000 ans.

Ce n’est pas sans émotion que nous nous sommes engagés sur la pente raide mais facile à
descendre, nous les premiers hommes… « Pense un peu, Frank, les premiers à pénétrer dans
ce mystère ! ! ! … ».

En bas se trouvait une vasque peu profonde, pleine d’une eau parfaitement limpide. Elle était
alimentée par une goutte d’eau qui tombait, à intervalles réguliers, par une sorte de cheminée
verticale qui s’ouvrait dans la voûte à moins d’un mètre du niveau de l’eau. Cette goutte
devait avoir une grande force car en tombant elle formait des éclaboussures et produisait une
note cristalline assez sonore pour être répercutée par la voûte de la grande salle. J’étais ravi ;
j’aurais voulu rester longtemps à l’écouter. Mais mon cousin voulait que nous continuions
l’exploration. Pendant plus d’une heure nous nous acharnâmes à chercher de nouvelles
galeries, mais nous n’avions pas de corde pour nous assurer et les quelques trous qu’on
pouvait discerner dans le rocher étaient trop haut pour qu’on puisse les atteindre avec nos
propres moyens.

Lorsque nous fûmes sortis Frank ne me cacha pas sa déception : \textit{« C’est vrai nous avons
trouvé une salle, me dit-il, mais elle n’a pas de stalagmites comme le clocher de Dargilan,
pas de stalactites excentriques comme à Fond de Gaume, pas de peintures rupestres comme à
Altamira, pas même un squelette d’ours ou de renard comme à Nabrigas, rien
d’intéressant… ».}
Le lendemain nous rendîmes la clé à son propriétaire lui racontant notre exploit ? \textit{« c’est ce
que je vous avais dit, concluait-il en haussant les épaules, il n’y a rien à tirer de cette
caverne, elle ne vaut même pas les frais d’une serrure ».}
Le soir un orage éclata suivi par une période de mauvais temps. C’était la fin des vacances et
des aventures
